
% Repository integration file.
% Assumptions:
%   - amsmath, amssymb, mathtools, amsthm are loaded in the parent document.
%   - theorem environments theorem, lemma, proposition, corollary, definition, conjecture,
%     remark, example are available.
% This file is intentionally written as a section/chapter fragment, not as a standalone document.

\providecommand{\MC}{\operatorname{MC}}
\providecommand{\Conf}{\operatorname{Conf}}
\providecommand{\Bar}{\operatorname{Bar}}
\providecommand{\Def}{\operatorname{Def}}
\providecommand{\Hom}{\operatorname{Hom}}
\providecommand{\im}{\operatorname{im}}
\providecommand{\id}{\operatorname{id}}
\providecommand{\op}{\mathrm{op}}
\providecommand{\gr}{\operatorname{gr}}
\providecommand{\pr}{\operatorname{pr}}
\providecommand{\Tot}{\operatorname{Tot}}
\providecommand{\Span}{\operatorname{Span}}
\providecommand{\Spec}{\operatorname{Spec}}
\providecommand{\Ann}{\operatorname{Ann}}
\providecommand{\Aut}{\operatorname{Aut}}
\providecommand{\Ord}{\mathrm{ord}}
\providecommand{\Coeff}{\mathrm{coeff}}
\providecommand{\Geom}{\mathrm{geom}}
\providecommand{\YB}{\mathrm{YB}}
\providecommand{\Ex}{\mathrm{Ex}}
\providecommand{\Ext}{\mathrm{ext}}
\providecommand{\full}{\mathrm{full}}
\providecommand{\pf}{\mathrm{pf}}
\providecommand{\conn}{\mathrm{conn}}
\providecommand{\ex}{\mathrm{ex}}
\providecommand{\bigtensor}{\mathbin{\widehat{\otimes}}}
\providecommand{\Bbbk}{\mathbb{k}}

\section{Modular dg-shifted Yangians from first principles}
\label{sec:modular-dg-shifted-yangians-first-principles}

\subsection{Orientation, status, and the Platonic target}

The genus-zero line-operator package and the modular bar package suggest one and the same
object viewed from two faces. On the one hand, perturbative line operators in a
three-dimensional holomorphic-topological theory are expected to be controlled by a Koszul-dual
$A_\infty$ algebra $A^!$ equipped with a universal Maurer--Cartan kernel
\[
r(z)\in A^!\bigtensor A^!((z^{-1}))
\]
satisfying an $A_\infty$ Yang--Baxter relation and a twisted coproduct.
On the other hand, the modular programme insists that higher genus is not external decoration
but a deformation variable internal to bar--cobar duality itself: the genus tower should be
governed by a universal Maurer--Cartan class
\[
\Theta_A \in \MC\!\left(\Def_{\mathrm{cyc}}(A)\bigtensor R\Gamma(\overline{\mathcal M}_{g,\bullet},\Bbbk)\right),
\]
whose scalar shadow is already visible in the Hodge bundle/Hodge class formalism and whose full
non-scalar existence remains the principal open problem of the modular homotopy programme.

The purpose of this section is to rebuild the ordered modular line-operator story from the ground
up, isolate the correct deformation object, and push the theory to a theorematic point where the
remaining frontier is reduced to an explicit cohomology problem. The guiding principles are:

\begin{enumerate}
\item
The modular $r$-kernel is \emph{not} the raw binary Taylor coefficient of $\Theta_A$.
It is the two-leg coefficient of the \emph{$\Theta_A$-deformed ordered twisting morphism}.
\item
Higher-point ordered corrections should be organized by \emph{exact support} rather than by
a loose ``number of external legs'' slogan.
\item
All higher exact-support corrections are governed by one pronilpotent dg Lie algebra.
\item
In the prime-form/Hodge realization, the first obstruction group $H^2$ is forced into the
three-leg sector; under a modular Arnold-type normalization, it becomes precisely the
modular Yang--Baxter class.
\end{enumerate}

The section is deliberately split into three layers.

\begin{enumerate}
\item[(I)]
\emph{Formal modular algebra.}
Here we define the ordered modular twisting system, the cumulant tower, and the exact-support
deformation dg Lie algebra. These constructions are new and theorematic once the ordered modular
bar package exists.

\item[(II)]
\emph{Prime-form/Hodge realization.}
Here we pass from formal algebra to a graph-decorated geometric model and prove an $H^2$-reduction
theorem. This is again theorematic, conditional only on a precise realization axiom.

\item[(III)]
\emph{The Platonic conjecture.}
Here we state the modular dg-shifted Yangian as the canonical completion of the genus-zero
line-operator Yangian package.
\end{enumerate}

\begin{remark}[What is input and what is new]
The following input is taken as established background.
\begin{enumerate}
\item
At genus zero, perturbative line operators are organized by a meromorphic tensor category
equivalent to modules over $A^!$, with OPE repackaged by a universal Maurer--Cartan kernel $r(z)$
and a twisted coproduct.
\item
In the bar--cobar language, the genus-zero $r(z)$ should be identified with the degree-$2$
piece of the twisting morphism, evaluated at the spectral parameter.
\item
In the modular programme, the scalar shadow of a universal class $\Theta_A$ is already visible,
while the full non-scalar $\Theta_A$ is still conjectural.
\end{enumerate}
Everything beyond those inputs---the ordered modular twisting system, the cumulant calculus, the
exact-support deformation dg Lie algebra, the Kuranishi theory, the planar-tree expansion, the
prime-form/Hodge $H^2$-reduction, and the modular Yang--Baxter identification theorem---is new
mathematical content developed here.
\end{remark}

\subsection{The ordered modular bar datum}

We start from the universal geometric datum that a higher-genus ordered OPE theory should use.

\begin{definition}[Ordered modular bar datum]
\label{def:ordered-modular-bar-datum}
An \emph{ordered modular bar datum} consists of the following data.

\begin{enumerate}
\item
A smooth curve $X$ and the universal stable ordered family
\[
\pi\colon \mathcal C \longrightarrow \mathcal M
\]
over an appropriate moduli stack $\mathcal M$ of smooth or stable pointed curves.

\item
For every finite ordered set $I$, an ordered compactified configuration space
\[
\overline{\Conf}^{\,\Ord}_I(\mathcal C/\mathcal M),
\]
together with restriction maps under order-preserving injections and forgetting maps.

\item
An augmented $E_1$-type object $A$ on $X$, together with a completed ordered bar coalgebra
\[
\widehat{\mathrm{Bar}}^{\,\Ord}_{\mathrm{mod}}(A)
=
\prod_{I}
\Omega^\bullet_{\log}\!\left(\overline{\Conf}^{\,\Ord}_I(\mathcal C/\mathcal M)\right)
\bigtensor
(s\bar A)^{\otimes I},
\]
equipped with the ordered deconcatenation coproduct.

\item
A prime-form/Hodge propagator system
\[
\eta^{\pf}_{ij}\in
\Omega^1_{\log}\!\left(\overline{\Conf}^{\,\Ord}_I(\mathcal C/\mathcal M),\mathcal K^{\pf}_{ij}\right)
\]
for every ordered pair of labels $i<j$, where $\mathcal K^{\pf}_{ij}$ is the corresponding
prime-form/Hodge coefficient line. The system is required to satisfy:
\begin{enumerate}
\item local singularity:
\[
\eta^{\pf}_{ij}\sim d\log(z_i-z_j)\qquad\text{near }z_i=z_j;
\]
\item antisymmetry:
\[
\eta^{\pf}_{ji}=-\eta^{\pf}_{ij};
\]
\item compatibility with clutching of curves;
\item compatibility with Verdier duality/opposition.
\end{enumerate}

\item
A completed ordered twisting morphism
\[
\tau\colon \widehat{\mathrm{Bar}}^{\,\Ord}_{\mathrm{mod}}(A)\to A^!
\]
solving the completed bar--cobar Maurer--Cartan equation
\[
\partial \tau+\tau\star\tau=0.
\]
\end{enumerate}
\end{definition}

\begin{remark}[The correct two-leg extraction]
\label{rem:correct-two-leg-extraction}
The modular two-leg kernel is not the binary Taylor coefficient of a deformation parameter.
It is the bar-length-$2$ coefficient of the \emph{deformed twisting morphism}. This is the ordered
higher-genus analogue of the genus-zero identification
\[
r(z)=\tau|_{\deg=2}.
\]
If $\Theta_A$ is the universal Maurer--Cartan class of the cyclic deformation complex, then the
modular two-leg extraction must be applied only after deforming $\tau$ by $\Theta_A$.
\end{remark}

\begin{definition}[Ordered modular two-leg extraction (conjectural)]
\label{defconj:ordered-modular-two-leg-extraction}
Assume the full non-scalar universal class
\[
\Theta_A\in
\MC\!\left(\Def_{\mathrm{cyc}}(A)\bigtensor R\Gamma(\overline{\mathcal M}_{g,\bullet},\Bbbk)\right)
\]
exists and deforms the ordered modular bar differential to a square-zero differential
$D_{\Theta_A}$ on $\widehat{\mathrm{Bar}}^{\,\Ord}_{\mathrm{mod}}(A)$.
Let
\[
\tau_{\Theta_A}\colon
\left(\widehat{\mathrm{Bar}}^{\,\Ord}_{\mathrm{mod}}(A),D_{\Theta_A}\right)\to A^!
\]
be the corresponding completed ordered twisting morphism.

Write
\[
\tau_{\Theta_A}=\sum_{m\ge 1}\tau_{\Theta_A}^{[m]}
\]
for its bar-length decomposition. Let
\[
\mu\in A^!\bigtensor A
\]
denote the universal coupling. Define the ordered two-leg extraction by
\[
\Ex^{(2)}_{\Ord}(\Theta_A)
:=
\operatorname{ev}_2\circ \operatorname{KD}_{\mu}\circ
\pr_{\mathrm{bar}=2}(\tau_{\Theta_A}),
\]
where $\operatorname{KD}_{\mu}$ contracts the two $A$-legs against two copies of $\mu$ and
$\operatorname{ev}_2$ evaluates on an ordered pair of marked points.

For every smooth pointed curve $(\Sigma;z,w)$ of genus $g$, define
\[
r_{g,\Sigma}(z,w)
:=
\Ex^{(2)}_{\Ord}(\Theta_A)\big|_{(\Sigma;z,w)}.
\]
The family
\[
r_A^{\mathrm{mod}}=\{r_{g,\Sigma}\}_{g,\Sigma}
\]
is the \emph{modular two-leg kernel} of $A$, and
\[
\mathsf Y^{\mathrm{mod}}(A):=(A^!,r_A^{\mathrm{mod}})
\]
is the \emph{modular dg-shifted Yangian} attached to $A$.
\end{definition}

\begin{remark}[The Platonic ideal]
The point of the definition-conjecture is that it isolates the \emph{correct object}.
Once the ordered modular bar package exists and the two-leg extraction is defined as above, all
higher compatibilities---clutching, Verdier duality, Yang--Baxter, higher ordered OPE coherence---are
shadows of a single Maurer--Cartan equation. The problem is thereby changed from
``guess higher-genus $r$-matrices'' to ``construct the ordered modular twisting system and solve its
exact-support deformation problem''.
\end{remark}

\subsection{Ordered modular twisting systems}

We now abstract the exact formal structure carried by the ordered modular bar package.

\begin{definition}[Ordered modular twisting system]
\label{def:ordered-modular-twisting-system}
An \emph{ordered modular twisting system} consists of a family
\[
\mathbb L=\{L_I\}_{I}
\]
indexed by finite ordered sets, together with the following data.

\begin{enumerate}
\item
For every ordered set $I$, a complete exact-support decomposition
\[
L_I=\prod_{\varnothing\neq S\subseteq I}L_I\langle S\rangle,
\]
where $L_I\langle S\rangle$ is the component supported exactly on the legs in $S$.

\item
For every order-preserving injection $f\colon J\hookrightarrow I$, an insertion map
\[
\iota_f\colon L_J\to L_I
\]
whose image lies in $L_I\langle f(J)\rangle$, functorial in $f$.

\item
For each ordered set $I$, an $L_\infty$-structure
\[
\ell_k^I\colon L_I^{\otimes k}\to L_I[2-k]
\]
satisfying support locality:
\[
\ell_k^I\big(L_I\langle S_1\rangle,\dots,L_I\langle S_k\rangle\big)
\subseteq
L_I\langle S_1\cup\cdots\cup S_k\rangle.
\]

\item
Connected locality: if the intersection graph on $\{1,\dots,k\}$ determined by
\[
a\sim b \quad\Longleftrightarrow\quad S_a\cap S_b\neq\varnothing
\]
is disconnected, then
\[
\ell_k^I(x_1,\dots,x_k)=0
\qquad
(x_a\in L_I\langle S_a\rangle).
\]

\item
For every separating clutching decomposition $I=I_-\sqcup I_+$, a bilinear node-contraction
\[
\star_\kappa\colon
L_{I_-\sqcup\{\bullet\}}\bigtensor L_{I_+\sqcup\{\circ\}}
\to
L_{I_-\sqcup I_+},
\]
such that external support is additive:
\[
\operatorname{supp}_{\Ext}(x\star_\kappa y)
=
\operatorname{supp}_{\Ext}(x)\sqcup \operatorname{supp}_{\Ext}(y).
\]

\item
A Verdier/opposition anti-involution
\[
\mathbb D_I\colon L_I\to L_{I^{\op}}
\]
compatible with support projections, insertion maps, and Maurer--Cartan formation.
\end{enumerate}
\end{definition}

\begin{remark}
In the ordered modular bar realization one should keep in mind
\[
L_I=\Hom^{\mathrm{cont}}(B_I(A),A^!),
\]
with exact support supplied by the support decomposition along ordered compactified boundary strata,
insertion by forgetting legs, clutching by nodal gluing with prime-form propagator, and Verdier
duality by duality of logarithmic forms.
\end{remark}

\begin{definition}[Pairwise seed, aggregate curvature, and cumulants]
\label{def:pairwise-seed-cumulants}
Let $\rho\in L_{\{1,2\}}^1$. For $i<j\in I$, write
\[
\rho_{ij}:=\iota_{\{1,2\}\hookrightarrow\{i,j\}}(\rho)\in L_I^1.
\]
Define the pairwise field on $I$ by
\[
R_I(\rho):=\sum_{i<j}\rho_{ij}.
\]
Its aggregate curvature is
\[
K_I(\rho)
:=
\MC_I\!\big(R_I(\rho)\big)
=
\sum_{k\ge 1}\frac1{k!}\,\ell_k^I\big(R_I(\rho),\dots,R_I(\rho)\big).
\]
Its \emph{ordered cumulant} is the full-support projection
\[
\mathfrak C_I(\rho):=\pr_I^{\full}\,K_I(\rho)\in L_I\langle I\rangle.
\]
\end{definition}

\begin{theorem}[Support decomposition of curvature]
\label{thm:support-decomposition-of-curvature}
For every finite ordered set $I$,
\[
K_I(\rho)
=
\sum_{\substack{J\subseteq I\\ |J|\ge 2}}
\iota_{J\hookrightarrow I}\,\mathfrak C_J(\rho).
\]
\end{theorem}

\begin{proof}
Decompose $K_I(\rho)$ into its exact-support pieces:
\[
K_I(\rho)=\sum_{\substack{J\subseteq I\\ |J|\ge 2}}K_{I;J}(\rho),
\qquad
K_{I;J}(\rho)\in L_I\langle J\rangle.
\]
By support locality, each summand in the expansion of $K_I(\rho)$ is supported on the union of
the supports of the participating pairwise legs, hence on some subset $J\subseteq I$.

Fix $J\subseteq I$. Forgetting all legs in $I\setminus J$ sends $R_I(\rho)$ to $R_J(\rho)$ and kills
every exact-support component not contained in $J$. Therefore the exact-support $J$-component
of $K_I(\rho)$ is precisely the inserted full-support component of $K_J(\rho)$:
\[
K_{I;J}(\rho)=\iota_{J\hookrightarrow I}\,\mathfrak C_J(\rho).
\]
Summing over all $J$ gives the result.
\end{proof}

\begin{theorem}[Boolean M\"obius inversion for cumulants]
\label{thm:boolean-mobius-inversion}
For every finite ordered set $I$,
\[
\mathfrak C_I(\rho)
=
\sum_{\substack{J\subseteq I\\ |J|\ge 2}}
(-1)^{|I|-|J|}
\,\iota_{J\hookrightarrow I}\,K_J(\rho).
\]
\end{theorem}

\begin{proof}
The previous theorem is exactly the zeta-transform relation
\[
K_I=\sum_{J\subseteq I}\mathfrak C_J^\uparrow
\]
on the Boolean lattice of subsets of $I$, where $\mathfrak C_J^\uparrow$ denotes insertion of
$\mathfrak C_J$ into $I$. The M\"obius function of the Boolean lattice is
\[
\mu(J,I)=(-1)^{|I|-|J|}.
\]
Applying M\"obius inversion yields the formula.
\end{proof}

\begin{corollary}[All-arity cumulant criterion]
\label{cor:all-arity-cumulant-criterion}
The pairwise field $R_I(\rho)$ is Maurer--Cartan for every ordered set $I$ if and only if
\[
\mathfrak C_I(\rho)=0
\qquad\text{for every }|I|\ge 2.
\]
\end{corollary}

\begin{proof}
If every cumulant vanishes, then by Theorem~\ref{thm:support-decomposition-of-curvature} every
aggregate curvature vanishes. Conversely, if $K_I(\rho)=0$, then
\[
\mathfrak C_I(\rho)=\pr_I^{\full}K_I(\rho)=0.
\]
\end{proof}

\begin{theorem}[Strict collapse to the two- and three-leg sectors]
\label{thm:strict-collapse-to-two-three}
Assume that each $L_I$ is a dg Lie algebra, so that only
\[
\ell_1=d,
\qquad
\ell_2=[-,-]
\]
are nonzero. Then
\[
\mathfrak C_I(\rho)=0
\qquad\text{for every }|I|\ge 4.
\]
Hence the pairwise field $R_I(\rho)$ is Maurer--Cartan for every ordered set $I$ if and only if
\[
\mathfrak C_{\{1,2\}}(\rho)=0
\qquad\text{and}\qquad
\mathfrak C_{\{1,2,3\}}(\rho)=0.
\]
\end{theorem}

\begin{proof}
A unary term $d(\rho_{ij})$ has support $\{i,j\}$, so it cannot contribute to full support for
$|I|\ge 3$.

A binary term $[\rho_{ij},\rho_{k\ell}]$ can be nonzero only when the supports intersect, by
connected locality. Therefore
\[
\{i,j\}\cap\{k,\ell\}\neq\varnothing,
\]
so the union support has size at most $3$. Hence no full-support contribution can occur on
$|I|\ge 4$.
The final statement follows from Corollary~\ref{cor:all-arity-cumulant-criterion}.
\end{proof}

\begin{theorem}[Clutching factorization of cumulants]
\label{thm:clutching-factorization-of-cumulants}
Assume that aggregate curvatures satisfy the clutching identity
\[
\xi^*K_{I_-\sqcup I_+}(\rho)
=
K_{I_-\sqcup\{\bullet\}}(\rho)\star_\kappa
K_{I_+\sqcup\{\circ\}}(\rho).
\]
Then the cumulants satisfy the exact factorization law
\[
\xi^*\mathfrak C_{I_-\sqcup I_+}(\rho)
=
\mathfrak C_{I_-\sqcup\{\bullet\}}(\rho)\star_\kappa
\mathfrak C_{I_+\sqcup\{\circ\}}(\rho).
\]
\end{theorem}

\begin{proof}
Expand the two factors on the right using Theorem~\ref{thm:support-decomposition-of-curvature}:
\[
K_{I_-\sqcup\{\bullet\}}
=
\sum_{P\subseteq I_-}
\iota_{P\sqcup\{\bullet\}}\mathfrak C_{P\sqcup\{\bullet\}},
\qquad
K_{I_+\sqcup\{\circ\}}
=
\sum_{Q\subseteq I_+}
\iota_{Q\sqcup\{\circ\}}\mathfrak C_{Q\sqcup\{\circ\}}.
\]
Hence
\[
\xi^*K_{I_-\sqcup I_+}
=
\sum_{P,Q}
\Bigl(
\mathfrak C_{P\sqcup\{\bullet\}}\star_\kappa
\mathfrak C_{Q\sqcup\{\circ\}}
\Bigr).
\]
By additivity of external support under $\star_\kappa$, the summand indexed by $(P,Q)$ has external
support exactly $P\sqcup Q$. Therefore the full-support projection to $I_-\sqcup I_+$ retains only the
term
\[
P=I_-,
\qquad
Q=I_+.
\]
This gives the desired identity.
\end{proof}

\begin{theorem}[Verdier transport of the cumulant tower]
\label{thm:verdier-transport-of-cumulants}
Assume $\mathbb D$ commutes with Maurer--Cartan formation and reverses ordered two-leg seeds:
\[
\mathbb D(\rho_{ij})=(\mathbb D\rho)_{ji}.
\]
Then
\[
\mathbb D K_I(\rho)=K_{I^{\op}}(\mathbb D\rho),
\qquad
\mathbb D\mathfrak C_I(\rho)=\mathfrak C_{I^{\op}}(\mathbb D\rho).
\]
\end{theorem}

\begin{proof}
By functoriality of insertion,
\[
\mathbb D\big(R_I(\rho)\big)=R_{I^{\op}}(\mathbb D\rho).
\]
Since $\mathbb D$ intertwines Maurer--Cartan formation,
\[
\mathbb D K_I(\rho)=K_{I^{\op}}(\mathbb D\rho).
\]
Projecting to full support yields the second identity.
\end{proof}

\subsection{The exact-support deformation dg Lie algebra}

The cumulant tower still treats each support separately. The next step is the decisive one:
all higher full-support corrections are controlled by a \emph{single} dg Lie algebra.

\begin{definition}[Pairwise Maurer--Cartan background]
\label{def:pairwise-mc-background}
Let $\rho\in L_{\{1,2\}}^1$ be a two-leg seed. We say that $\rho$ is a
\emph{pairwise Maurer--Cartan background} if
\[
K_I(\rho)=0
\qquad\text{for every ordered set }I.
\]
By Theorem~\ref{thm:strict-collapse-to-two-three}, in a strict dg Lie system it suffices to check the
two-leg Maurer--Cartan equation and the three-leg Yang--Baxter relation.
\end{definition}

Fix such a background $\rho$ for the rest of this subsection. For each ordered set $I$, write
\[
R_I:=R_I(\rho).
\]

\begin{definition}[Exact-support deformation dg Lie algebra]
\label{def:exact-support-deformation-dg-Lie-algebra}
Define
\[
\mathfrak g_{\ex}(\rho)
:=
\prod_{|I|\ge 3}L_I\langle I\rangle.
\]
For $x=\{x_I\}\in \mathfrak g_{\ex}(\rho)$, define its insertion envelope on $I$ by
\[
\widetilde x_I
:=
\sum_{\substack{J\subseteq I\\ |J|\ge 3}}
\iota_{J\hookrightarrow I}(x_J).
\]

Define
\[
(\delta_\rho x)_I
:=
\pr_I^{\full}\bigl(d\widetilde x_I+[R_I,\widetilde x_I]\bigr),
\]
and
\[
[x,y]_{\ex,I}
:=
\pr_I^{\full}[\widetilde x_I,\widetilde y_I].
\]
\end{definition}

\begin{lemma}[Support-envelope descent]
\label{lem:support-envelope-descent}
For every $x,y\in \mathfrak g_{\ex}(\rho)$ and every ordered set $I$,
\[
d\widetilde x_I+[R_I,\widetilde x_I]
=
\widetilde{(\delta_\rho x)}_I,
\]
and
\[
[\widetilde x_I,\widetilde y_I]
=
\widetilde{[x,y]_{\ex}}_I.
\]
\end{lemma}

\begin{proof}
Consider the exact-support decomposition of the first expression. Because $d$ preserves support and
the bracket is support-local, the exact-support $S$-part of
\[
d\widetilde x_I+[R_I,\widetilde x_I]
\]
depends only on the pieces of $\widetilde x_I$ supported on subsets of $S$ together with the pairwise
edges inside $S$. By functoriality of insertion, that exact-support $S$-part is precisely the inserted copy of
\[
(\delta_\rho x)_S
=
\pr_S^{\full}\bigl(d\widetilde x_S+[R_S,\widetilde x_S]\bigr).
\]
Summing over all $S\subseteq I$ gives the first identity.

The proof of the second identity is identical:
the exact-support $S$-part of
\[
[\widetilde x_I,\widetilde y_I]
\]
comes only from brackets between inserted pieces supported on subsets of $S$, hence is the inserted copy
of
\[
[x,y]_{\ex,S}
=
\pr_S^{\full}[\widetilde x_S,\widetilde y_S].
\]
\end{proof}

\begin{theorem}[Exact-support dg Lie algebra]
\label{thm:exact-support-dg-Lie-algebra}
The triple
\[
\bigl(\mathfrak g_{\ex}(\rho),\delta_\rho,[\, ,\,]_{\ex}\bigr)
\]
is a complete dg Lie algebra.
\end{theorem}

\begin{proof}
Define
\[
D_I:=d+[R_I,-]
\]
on $L_I$. Since $R_I$ is Maurer--Cartan,
\[
D_I^2=[dR_I+\tfrac12[R_I,R_I],-]=0.
\]

By Lemma~\ref{lem:support-envelope-descent},
\[
\widetilde{(\delta_\rho^2 x)}_I
=
D_I\widetilde{(\delta_\rho x)}_I
=
D_I^2\widetilde x_I
=
0.
\]
The insertion-envelope functor is faithful on exact-support families:
if $\widetilde z_I=0$ for all $I$, then $z=0$, because choosing $I$ minimal with $z_I\neq 0$ yields
$\widetilde z_I=z_I$. Hence $\delta_\rho^2=0$.

For the derivation property,
\[
\widetilde{\delta_\rho[x,y]_{\ex}}_I
=
D_I[\widetilde x_I,\widetilde y_I]
=
[D_I\widetilde x_I,\widetilde y_I]+(-1)^{|x|}[\widetilde x_I,D_I\widetilde y_I]
=
\widetilde{[\delta_\rho x,y]_{\ex}+(-1)^{|x|}[x,\delta_\rho y]_{\ex}}_I,
\]
and faithfulness again gives the desired identity.

Finally,
\[
\widetilde{[x,[y,z]_{\ex}]_{\ex}}_I
=
[\widetilde x_I,[\widetilde y_I,\widetilde z_I]]
\]
and the Jacobi identity in $L_I$ descends to the exact-support bracket by faithfulness.
Completeness is inherited from the completed topology of the exact-support summands.
\end{proof}

\begin{theorem}[Equivalence with modular completions]
\label{thm:equivalence-with-modular-completions}
For $x\in\mathfrak g_{\ex}(\rho)^1$, define the total ordered field
\[
\mathcal A_I(x):=R_I+\widetilde x_I.
\]
Then the following are equivalent.

\begin{enumerate}
\item
$x$ is Maurer--Cartan in $\mathfrak g_{\ex}(\rho)$:
\[
\delta_\rho x+\frac12[x,x]_{\ex}=0.
\]

\item
For every ordered set $I$,
\[
d\mathcal A_I(x)+\frac12[\mathcal A_I(x),\mathcal A_I(x)]=0.
\]
\end{enumerate}
\end{theorem}

\begin{proof}
Since $R_I$ is Maurer--Cartan,
\[
d\mathcal A_I(x)+\frac12[\mathcal A_I(x),\mathcal A_I(x)]
=
D_I\widetilde x_I+\frac12[\widetilde x_I,\widetilde x_I].
\]
By Lemma~\ref{lem:support-envelope-descent},
\[
D_I\widetilde x_I+\frac12[\widetilde x_I,\widetilde x_I]
=
\widetilde{\delta_\rho x}_I+\frac12\widetilde{[x,x]_{\ex}}_I
=
\widetilde{\delta_\rho x+\frac12[x,x]_{\ex}}_I.
\]
Faithfulness of the insertion envelope implies the equivalence.
\end{proof}

\begin{remark}[Interpretation]
The exact-support dg Lie algebra is the correct deformation object. A modular completion of the
pairwise background $R_I$ is not a scattered family of higher-point corrections; it is a single Maurer--Cartan
element of $\mathfrak g_{\ex}(\rho)$.
\end{remark}

\subsection{Kuranishi theory and explicit tree formulas}

We now pass from abstract deformation theory to explicit construction.

Choose a contraction of cochain complexes
\[
(H \xrightarrow{i} \mathfrak g_{\ex}(\rho)\xrightarrow{p} H,\; h)
\]
satisfying
\[
\delta_\rho h+h\delta_\rho=\id-ip,
\qquad
pi=\id,
\qquad
ph=hi=h^2=0.
\]
Think of $H$ as the exact-support cohomology controlling modular periods.

\begin{theorem}[Kuranishi theorem for ordered modular completions]
\label{thm:kuranishi-theorem-for-ordered-modular-completions}
For every formal parameter $\kappa\in H^1$, the fixed-point equation
\[
x
=
i\kappa-\frac12\,h[x,x]_{\ex}
\tag{$\mathrm K$}
\]
admits a unique formal solution $x(\kappa)$ satisfying
\[
px(\kappa)=\kappa,
\qquad
hx(\kappa)=0.
\]

Define the Kuranishi map
\[
\mathcal K(\kappa):=\frac12\,p[x(\kappa),x(\kappa)]_{\ex}\in H^2.
\]
Then
\[
x(\kappa)\text{ is Maurer--Cartan}
\qquad\Longleftrightarrow\qquad
\mathcal K(\kappa)=0.
\]
\end{theorem}

\begin{proof}
Existence and uniqueness of the formal fixed point follow from pronilpotence:
starting from $x_0=i\kappa$ and iterating
\[
x_{n+1}=i\kappa-\frac12 h[x_n,x_n]_{\ex},
\]
each iteration gains filtration and converges formally to a unique solution.

Set
\[
F:=\delta_\rho x+\frac12[x,x]_{\ex}.
\]
From the fixed-point equation,
\[
x-i\kappa=-\frac12 h[x,x]_{\ex}.
\]
Applying $\delta_\rho$ and using
\[
\delta_\rho h=\id-ip-h\delta_\rho,
\]
one obtains
\[
F=ipF+h[F,x]_{\ex}.
\tag{$*$}
\]
Taking $p$ of both sides gives
\[
pF=\frac12\,p[x,x]_{\ex}=\mathcal K(\kappa).
\]
If $\mathcal K(\kappa)=0$, then $(*)$ becomes
\[
F=h[F,x]_{\ex}.
\]
Because the operator $F\mapsto h[F,x]_{\ex}$ raises filtration, pronilpotence forces $F=0$.
Thus $x(\kappa)$ is Maurer--Cartan. The converse is immediate.
\end{proof}

\begin{corollary}[Unobstructed regime]
\label{cor:unobstructed-regime}
If
\[
H^2\bigl(\mathfrak g_{\ex}(\rho)\bigr)=0,
\]
then every class $\kappa\in H^1(\mathfrak g_{\ex}(\rho))$ integrates to a gauge-normalized
modular completion. If moreover
\[
H^1\bigl(\mathfrak g_{\ex}(\rho)\bigr)=0,
\]
then the modular completion is unique up to gauge.
\end{corollary}

\begin{proof}
If $H^2=0$, the Kuranishi map vanishes identically, and
Theorem~\ref{thm:kuranishi-theorem-for-ordered-modular-completions} gives a Maurer--Cartan solution
for every $\kappa$. If also $H^1=0$, there is only the zero parameter, hence only one gauge-normalized
solution.
\end{proof}

\begin{definition}[Planar binary tree operators]
\label{def:planar-binary-tree-operators}
For a rooted planar binary tree $T$, define $\operatorname{Op}_T(\kappa)\in\mathfrak g_{\ex}(\rho)^1$
recursively by
\[
\operatorname{Op}_{\bullet}(\kappa):=i\kappa,
\]
and
\[
\operatorname{Op}_{T_1\vee T_2}(\kappa)
:=
h\bigl[\operatorname{Op}_{T_1}(\kappa),\operatorname{Op}_{T_2}(\kappa)\bigr]_{\ex}.
\]
\end{definition}

\begin{theorem}[Canonical planar-tree expansion]
\label{thm:canonical-planar-tree-expansion}
Assume $H^2(\mathfrak g_{\ex}(\rho))=0$.
Then the unique gauge-normalized Maurer--Cartan solution with cohomology parameter
$\kappa\in H^1$ is
\[
x(\kappa)
=
\sum_{T\in \mathrm{PBT}}
\left(-\frac12\right)^{|V_{\mathrm{int}}(T)|}
\operatorname{Op}_T(\kappa),
\]
where the sum ranges over rooted planar binary trees.
\end{theorem}

\begin{proof}
Iterating the Kuranishi equation
\[
x=i\kappa-\frac12 h[x,x]_{\ex}
\]
substitutes one new binary bracket and one new homotopy $h$ at each stage. Therefore every fully expanded
term is indexed by a rooted planar binary tree. Each internal vertex contributes a factor $-\tfrac12$.
Pronilpotence ensures formal convergence.
\end{proof}

\begin{remark}[Meaning of the tree formula]
The formula is not ornamental. It says that once one knows the exact-support bracket, the homotopy $h$,
and the tangent class $\kappa$, \emph{all higher modular corrections are explicit tree sums}.
If $h$ is realized by a prime-form/Hodge Green operator, this becomes an actual geometric construction,
not merely a formal existence statement.
\end{remark}

\subsection{Prime-form/Hodge graph realization}

We now pass to the realization best adapted to actual configuration-space geometry.

\begin{definition}[Prime-form/Hodge graph realization]
\label{def:prime-form-hodge-graph-realization}
A \emph{prime-form/Hodge graph realization} of the exact-support deformation dg Lie algebra
$\mathfrak g_{\ex}(\rho)$ consists of the following additional structure.

\begin{enumerate}
\item
For each ordered set $I$, the full-support summand $L_I\langle I\rangle$ is realized as the completion
of a graph-decorated complex
\[
\mathcal G_I^{\pf}(A^!)
=
\prod_{\Gamma\in \mathsf{Conn}(I)}
\Omega^\bullet_{\log}\!\bigl(\overline{\Conf}^{\,\Ord}_I(\mathcal C/\mathcal M),
\mathcal K_\Gamma^{\pf}\bigr)
\bigtensor
\mathcal C_\Gamma(A^!),
\]
where:
\begin{enumerate}
\item
$\mathsf{Conn}(I)$ is the set of connected graphs on vertex set $I$;
\item
$\mathcal K_\Gamma^{\pf}$ is the tensor product of prime-form/Hodge kernel lines attached to the edges of
$\Gamma$;
\item
$\mathcal C_\Gamma(A^!)$ is the coefficient complex attached to the vertex decorations and algebraic
operations in $A^!$.
\end{enumerate}

\item
Every edge of $\Gamma$ contributes a prime-form/Hodge propagator, hence a logarithmic one-form.
Therefore the \emph{geometric degree} of a graph term is exactly the number of edges:
\[
\deg_{\Geom}(\Gamma)=|E(\Gamma)|.
\]

\item
The exact-support differential decomposes as
\[
\delta_\rho
=
d_{\Coeff}+d_{\Geom}+\partial_R,
\]
where:
\begin{enumerate}
\item
$d_{\Coeff}$ preserves geometric degree;
\item
$d_{\Geom}$ raises geometric degree by $1$;
\item
$\partial_R$ also raises geometric degree by $1$.
\end{enumerate}

\item
The coefficient complexes are \emph{connective}:
\[
H^q\bigl(\mathcal C_\Gamma(A^!)\bigr)=0
\qquad\text{for }q<0.
\]
\end{enumerate}
\end{definition}

\begin{remark}[Why this is the right realization]
Connected full-support graphs are the natural expansion variables because exact support forces connectedness,
and every propagator edge carries a prime-form/Hodge $1$-form. The connectiveness assumption is the
geometric reflection of the usual deformation-theoretic principle that negative-degree coefficient
cohomology should not appear in the first obstruction range.
\end{remark}

\begin{lemma}[Connected graph edge bound]
\label{lem:connected-graph-edge-bound}
If $\Gamma$ is a connected graph on a finite set $I$, then
\[
|E(\Gamma)|\ge |I|-1.
\]
Consequently every full-support graph term on $I$ has geometric degree at least $|I|-1$.
\end{lemma}

\begin{proof}
Every connected graph on $|I|$ vertices contains a spanning tree, and every spanning tree has
$|I|-1$ edges.
\end{proof}

\begin{theorem}[Geometric-degree spectral sequence]
\label{thm:geometric-degree-spectral-sequence}
Let $F^p\mathfrak g_{\ex}(\rho)$ be the filtration by geometric degree:
\[
F^p\mathfrak g_{\ex}(\rho)
=
\{x\in \mathfrak g_{\ex}(\rho)\mid \deg_{\Geom}(x)\ge p\}.
\]
Then $F^\bullet$ is a complete descending filtration preserved by $\delta_\rho$, and it defines
a convergent spectral sequence
\[
E_1^{p,q}
\Longrightarrow
H^{p+q}\bigl(\mathfrak g_{\ex}(\rho)\bigr),
\]
whose $E_1$-page is
\[
E_1^{p,q}
\cong
\bigoplus_{|I|\ge 3}
H^q\!\left(\gr_F^p L_I\langle I\rangle,d_{\Coeff}\right).
\]
Moreover
\[
E_1^{p,q}=0
\qquad\text{unless }p\ge |I|-1.
\]
\end{theorem}

\begin{proof}
Because $d_{\Coeff}$ preserves geometric degree and $d_{\Geom},\partial_R$ raise it by $1$,
the filtration is preserved:
\[
\delta_\rho(F^p)\subseteq F^p.
\]
Completeness follows from completion in the graph-expansion topology.

The associated graded differential is precisely $d_{\Coeff}$, because the other pieces strictly increase
geometric degree. This gives the displayed form of the $E_1$-page. The vanishing statement follows from
Lemma~\ref{lem:connected-graph-edge-bound}.
\end{proof}

\begin{theorem}[Three-leg reduction of the first obstruction]
\label{thm:three-leg-reduction-of-H2}
Assume a prime-form/Hodge graph realization.
Then the natural inclusion of the full-support three-leg subcomplex induces an isomorphism
\[
H^2\bigl(\mathfrak g_{\ex}(\rho)\bigr)
\cong
H^2\bigl(L_{\{1,2,3\}}\langle \{1,2,3\}\rangle,\delta_\rho\bigr).
\]
Equivalently, the first obstruction is purely three-legged.
\end{theorem}

\begin{proof}
On the $E_1$-page of the spectral sequence of
Theorem~\ref{thm:geometric-degree-spectral-sequence}, total degree $2$ can only come from terms with
\[
p+q=2.
\]
By connectiveness of the coefficient complexes, $q<0$ is impossible. By
Lemma~\ref{lem:connected-graph-edge-bound}, the condition $p\ge |I|-1$ forces
\[
|I|-1\le 2,
\qquad |I|\ge 3,
\]
hence $|I|=3$ and $p=2$, $q=0$. Therefore the only possible contribution to total degree $2$
comes from the full-support three-leg sector.

No differential can either enter or leave that bidegree:
incoming differentials would require negative coefficient degree, outgoing differentials would land in
total degree $3$. Hence
\[
E_1^{2,0}=E_\infty^{2,0},
\]
and this identifies $H^2(\mathfrak g_{\ex}(\rho))$ with the three-leg full-support cohomology.
\end{proof}

\begin{remark}[Frontier shift]
This theorem is the first genuine reduction theorem of the modular line-operator programme.
It says that the existence of a modular dg-shifted Yangian does not depend on a mysterious infinite
tower of unrelated higher-genus equations. The first obstruction is already and only the three-leg
full-support class.
\end{remark}

\subsection{The three-leg sector and the modular Yang--Baxter class}

We now identify the three-leg obstruction more sharply.

\begin{definition}[Three-leg Y-graph complex]
\label{def:three-leg-Y-graph-complex}
Let
\[
I_3:=\{1,2,3\}.
\]
In the prime-form/Hodge graph realization, define the \emph{three-leg Y-graph subcomplex}
$\mathcal Y_\rho\subset L_{I_3}\langle I_3\rangle$ to be the geometric-degree-$2$ summand spanned by
the three connected spanning trees
\[
(12,13),\qquad (12,23),\qquad (13,23),
\]
with their coefficient decorations.
Concretely,
\[
\mathcal Y_\rho
=
\Span\Bigl\{
\eta_{12}^{\pf}\wedge \eta_{13}^{\pf}\otimes c_{12,13},\ 
\eta_{12}^{\pf}\wedge \eta_{23}^{\pf}\otimes c_{12,23},\ 
\eta_{13}^{\pf}\wedge \eta_{23}^{\pf}\otimes c_{13,23}
\Bigr\},
\]
where the $c_{ab,cd}$ lie in the relevant degree-$0$ coefficient spaces.
\end{definition}

\begin{definition}[Strict coefficient Yang--Baxter element]
\label{def:strict-coefficient-yb-element}
Assume the coefficient algebra is strict dg Lie in the three-leg sector.
Then the \emph{coefficient Yang--Baxter element} is
\[
\YB(\rho)
:=
[\rho_{12},\rho_{13}]
+
[\rho_{12},\rho_{23}]
+
[\rho_{13},\rho_{23}].
\]
\end{definition}

\begin{definition}[Modular Arnold normalization]
\label{def:modular-arnold-normalization}
A prime-form/Hodge propagator system is said to satisfy \emph{modular Arnold normalization} on the
ordered three-point sector if there exist:
\begin{enumerate}
\item a closed full-support $2$-form
\[
\Omega_{123}\in
\Omega^2_{\log}\!\bigl(\overline{\Conf}^{\,\Ord}_{I_3}(\mathcal C/\mathcal M),
\mathcal K^{\pf}_{123}\bigr),
\]
\item $1$-forms
\[
\beta_{12,13},\qquad \beta_{12,23},\qquad \beta_{13,23},
\]
such that
\[
\eta_{12}^{\pf}\wedge \eta_{13}^{\pf}
=
\Omega_{123}+d\beta_{12,13},
\]
\[
\eta_{12}^{\pf}\wedge \eta_{23}^{\pf}
=
\Omega_{123}+d\beta_{12,23},
\]
\[
\eta_{13}^{\pf}\wedge \eta_{23}^{\pf}
=
\Omega_{123}+d\beta_{13,23}.
\]
\end{enumerate}
\end{definition}

\begin{remark}[Meaning of modular Arnold normalization]
At genus zero, with $\eta_{ij}=d\log(z_i-z_j)$, the wedge products of the three propagators are all
rational multiples of the same ordered volume form, and the classical Arnold relation is the resulting
linear dependence. The definition above is the higher-genus prime-form/Hodge upgrade of that fact:
it asserts that the three Y-graph $2$-forms represent one common cohomology class in the full-support
three-point sector.
\end{remark}

\begin{theorem}[Modular Yang--Baxter identification of $H^2$]
\label{thm:modular-yang-baxter-identification-of-H2}
Assume:
\begin{enumerate}
\item
a prime-form/Hodge graph realization;
\item
strict dg Lie coefficients in the three-leg sector;
\item
modular Arnold normalization.
\end{enumerate}
Then the first obstruction group is canonically identified with the three-leg Yang--Baxter class:
\[
H^2\bigl(\mathfrak g_{\ex}(\rho)\bigr)
\cong
[\Omega_{123}]\,\cdot\,
H^0\!\bigl(\Span\{\YB(\rho)\},d_{\Coeff}\bigr).
\]
More concretely, the class of the full-support three-leg curvature equals
\[
\bigl[\Omega_{123}\bigr]\otimes \bigl[\YB(\rho)\bigr].
\]
\end{theorem}

\begin{proof}
By Theorem~\ref{thm:three-leg-reduction-of-H2}, the group
$H^2(\mathfrak g_{\ex}(\rho))$ is the three-leg full-support cohomology.
In geometric degree $2$, the only connected full-support graphs are the three spanning trees, so the
degree-$2$ curvature class has the form
\[
\eta_{12}^{\pf}\wedge\eta_{13}^{\pf}\otimes[\rho_{12},\rho_{13}]
+
\eta_{12}^{\pf}\wedge\eta_{23}^{\pf}\otimes[\rho_{12},\rho_{23}]
+
\eta_{13}^{\pf}\wedge\eta_{23}^{\pf}\otimes[\rho_{13},\rho_{23}].
\]
By modular Arnold normalization, each geometric factor is cohomologous to $\Omega_{123}$.
Therefore the cohomology class of the entire expression is
\[
[\Omega_{123}]
\otimes
\left(
[\rho_{12},\rho_{13}]
+
[\rho_{12},\rho_{23}]
+
[\rho_{13},\rho_{23}]
\right),
\]
namely
\[
[\Omega_{123}]\otimes [\YB(\rho)].
\]
This identifies the first obstruction group with the coefficient Yang--Baxter class multiplied by the
canonical geometric three-leg class.
\end{proof}

\begin{corollary}[Vanishing criterion for the first obstruction]
\label{cor:vanishing-criterion-first-obstruction}
Under the assumptions of
Theorem~\ref{thm:modular-yang-baxter-identification-of-H2}, one has
\[
H^2\bigl(\mathfrak g_{\ex}(\rho)\bigr)=0
\qquad\Longleftrightarrow\qquad
[\YB(\rho)]=0
\]
provided $[\Omega_{123}]\neq 0$.
In particular, if the coefficient Yang--Baxter element $\YB(\rho)$ is $d_{\Coeff}$-exact or vanishes,
then the modular completion problem is unobstructed in degree $2$.
\end{corollary}

\begin{proof}
Immediate from Theorem~\ref{thm:modular-yang-baxter-identification-of-H2}.
\end{proof}

\begin{remark}[The first real frontier theorem]
This is the theorem that pushes the programme forward in a decisive way.
It converts the existence problem for modular dg-shifted Yangians from a vague higher-genus aspiration
into an explicit statement:

\medskip

\centerline{\emph{the first obstruction is exactly the modular Yang--Baxter class on three legs}.}

\medskip

Everything higher begins only after the three-leg class has been killed.
\end{remark}

\subsection{A theorematic existence criterion}

We can now formulate the cleanest existence theorem available at this stage.

\begin{theorem}[Formal existence of modular dg-shifted Yangians]
\label{thm:formal-existence-of-modular-dg-shifted-yangians}
Assume:
\begin{enumerate}
\item
an ordered modular twisting system extracted from a completed ordered modular bar datum;
\item
a pairwise Maurer--Cartan background $\rho$;
\item
a prime-form/Hodge graph realization;
\item
modular Arnold normalization;
\item
vanishing of the first obstruction class:
\[
[\YB(\rho)]=0;
\]
\item
vanishing of higher obstruction groups:
\[
H^2\bigl(\gr_F^m\mathfrak g_{\ex}(\rho)\bigr)=0
\qquad\text{for all }m\ge 3.
\]
\end{enumerate}
Then there exists a formal modular completion
\[
x\in \MC\bigl(\mathfrak g_{\ex}(\rho)\bigr)
\]
and hence a compatible family of ordered higher-genus OPE fields
\[
\mathcal A_I(x)=R_I+\widetilde x_I
\]
solving the full ordered Maurer--Cartan system on every ordered sector $I$.

If in addition
\[
H^1\bigl(\gr_F^m\mathfrak g_{\ex}(\rho)\bigr)=0
\qquad\text{for all }m\ge 3,
\]
then the completion is unique up to gauge.
\end{theorem}

\begin{proof}
By Corollary~\ref{cor:vanishing-criterion-first-obstruction}, the first obstruction vanishes.
The higher obstruction groups vanish by assumption. Standard filtered dg Lie induction then lifts a
solution order by order in the geometric/support filtration. The resulting Maurer--Cartan element
exists in the completed dg Lie algebra $\mathfrak g_{\ex}(\rho)$.
Uniqueness follows from the vanishing of the corresponding $H^1$-groups.
\end{proof}

\begin{corollary}[Explicit tree formula in the unobstructed case]
\label{cor:explicit-tree-formula-unobstructed}
If
\[
H^2\bigl(\mathfrak g_{\ex}(\rho)\bigr)=0,
\]
then every gauge-normalized modular completion is given by the planar-tree expansion of
Theorem~\ref{thm:canonical-planar-tree-expansion}.
\end{corollary}

\subsection{The formal definition--conjecture to carry into the repository}

The previous subsections separate the theorematic heart from the still-conjectural universal input.
We now package the whole story in the form that should sit in the research repository.

\begin{definition}[Modular dg-shifted Yangian of an ordered modular Koszul object (conjectural)]
\label{defconj:modular-dg-shifted-yangian}
Let $A$ be a modular Koszul $E_1$-type object on a smooth curve $X$, with Koszul dual $A^!$.
Assume:

\begin{enumerate}
\item[(H1)]
the ordered modular bar datum of Definition~\ref{def:ordered-modular-bar-datum} exists;

\item[(H2)]
the cyclic deformation complex $\Def_{\mathrm{cyc}}(A)$ exists as a complete cyclic $L_\infty$-algebra,
together with a universal Maurer--Cartan class
\[
\Theta_A\in
\MC\!\left(\Def_{\mathrm{cyc}}(A)\bigtensor R\Gamma(\overline{\mathcal M}_{g,\bullet},\Bbbk)\right)
\]
compatible with clutching and Verdier duality;

\item[(H3)]
the $\Theta_A$-deformed ordered twisting morphism $\tau_{\Theta_A}$ exists;

\item[(H4)]
the genus-zero fiber recovers the line-operator dg-shifted Yangian package:
\[
r_{0,\mathbb{P}^1}(z,w)=r(z-w).
\]
\end{enumerate}

Then the \emph{modular dg-shifted Yangian} of $A$ is the pair
\[
\mathsf Y^{\mathrm{mod}}(A)
=
(A^!,r_A^{\mathrm{mod}}),
\qquad
r_A^{\mathrm{mod}}:=\Ex_{\Ord}^{(2)}(\Theta_A).
\]

Conjecturally, its completed smooth module category is the modular meromorphic tensor category of
perturbative line operators associated to $A$.

This object should be characterized by the following formal properties.

\begin{enumerate}
\item[(P1)]
\textbf{Genus-zero normalization and local singularity.}
Near the diagonal,
\[
r_{g,\Sigma}(z,w)=r_{\mathrm{sing}}(z,w)+r^{\mathrm{reg}}_{g,\Sigma}(z,w),
\]
where $r_{\mathrm{sing}}$ is the universal genus-zero singular piece and
\[
r_{0,\mathbb{P}^1}(z,w)=r(z-w).
\]

\item[(P2)]
\textbf{Clutching/factorization.}
For every boundary clutching morphism,
\[
\xi^*r_A^{\mathrm{mod}}
=
r_{A,-}^{\mathrm{mod}}\star_\kappa r_{A,+}^{\mathrm{mod}}.
\]

\item[(P3)]
\textbf{Verdier/opposition.}
Under Verdier--Koszul duality,
\[
\mathbb D\bigl(r_A^{\mathrm{mod}}(z,w)\bigr)=\bigl(r_{A^!}^{\mathrm{mod}}(w,z)\bigr)^{21}.
\]

\item[(P4)]
\textbf{Modular Yang--Baxter/sewing law.}
The three-leg exact-support component of the ordered Maurer--Cartan equation is the modular
Yang--Baxter relation. In the strict coefficient regime it reduces to the coefficient identity
\[
[\rho_{12},\rho_{13}]
+
[\rho_{12},\rho_{23}]
+
[\rho_{13},\rho_{23}]
=
0.
\]
\end{enumerate}
\end{definition}

\subsection{What has actually been proved in this section}

For future use and for repository clarity, we summarize the theorematic content in one place.

\begin{theorem}[The theorematic core]
\label{thm:theorematic-core}
Assume an ordered modular twisting system and a pairwise Maurer--Cartan background $\rho$.
Then the following hold.

\begin{enumerate}
\item
The full ordered curvature decomposes into exact-support cumulants
(Theorem~\ref{thm:support-decomposition-of-curvature}), and cumulants are recovered from
aggregate curvatures by Boolean M\"obius inversion
(Theorem~\ref{thm:boolean-mobius-inversion}).

\item
In a strict dg Lie system, all higher cumulants vanish automatically for $|I|\ge 4$; hence
the full pairwise Maurer--Cartan problem collapses to the two- and three-leg sectors
(Theorem~\ref{thm:strict-collapse-to-two-three}).

\item
The family of exact-support corrections is governed by one complete dg Lie algebra
\[
\mathfrak g_{\ex}(\rho)
\]
(Theorem~\ref{thm:exact-support-dg-Lie-algebra}), and Maurer--Cartan elements of this dg Lie algebra
are equivalent to full ordered modular completions (Theorem~\ref{thm:equivalence-with-modular-completions}).

\item
The formal moduli of modular completions is controlled by the Kuranishi map
\[
\mathcal K\colon H^1\to H^2
\]
(Theorem~\ref{thm:kuranishi-theorem-for-ordered-modular-completions}), and unobstructed completions
admit an explicit planar-tree expansion
(Theorem~\ref{thm:canonical-planar-tree-expansion}).

\item
In the prime-form/Hodge graph realization, the first obstruction group is purely three-legged
(Theorem~\ref{thm:three-leg-reduction-of-H2}).

\item
Under modular Arnold normalization, the first obstruction is canonically identified with the
coefficient Yang--Baxter class
(Theorem~\ref{thm:modular-yang-baxter-identification-of-H2}).
\end{enumerate}
\end{theorem}

\subsection{Research consequences and the next frontier}

The formalism built here shifts the frontier in a decisive way.

\begin{enumerate}
\item
The modular line-operator problem is no longer the search for ad hoc higher-genus kernels.
It is the Maurer--Cartan problem of a single dg Lie algebra,
\[
\mathfrak g_{\ex}(\rho).
\]

\item
The obstruction theory is no longer dispersed over all arities.
Its first nontrivial term is canonically three-legged and geometric.

\item
The prime-form/Hodge realization makes the first obstruction explicitly calculable:
one computes the Y-graph class on the ordered three-point compactification.

\item
The most concrete next theorem is now sharply posed:

\begin{center}
\emph{identify or prove the vanishing of }
\[
H^2\bigl(\mathfrak g_{\ex}(\rho)\bigr)
\]
\emph{for the prime-form/Hodge realization of the ordered modular bar complex.}
\end{center}

By Theorem~\ref{thm:modular-yang-baxter-identification-of-H2}, this reduces in the strict regime to
the modular Yang--Baxter class.
\end{enumerate}

\begin{remark}[The Platonic ideal, restated]
The final philosophical synthesis is simple.

\begin{center}
\emph{Genus zero gives the seed, exact support gives the deformation object,}\\
\emph{prime-form/Hodge geometry gives the obstruction class, and}\\
\emph{modular homotopy should produce the universal Maurer--Cartan lift.}
\end{center}

This is the ordered extended-operator face of modular homotopy.
\end{remark}
